% future/cpu.tex
% mainfile: ../perfbook.tex
% SPDX-License-Identifier: CC-BY-SA-3.0

\section{CPU技术的未来已今非昔比}
\label{sec:future:The Future of CPU Technology Ain't What it Used to Be}
%
\epigraph{辉煌的未来已成过去。}{David Maraniss}

根据多年的经验，回首过去的岁月总是那么简单而纯真。
21世纪初，最大的无知表现在：大多数人对\IXr{Moore's Law}
即将无法继续带来传统CPU时钟频率提升这一事实浑然不知。
当然，偶尔也有关于技术极限的警告，但这样的警告已经发出了数十年。
怀着这样的心态，请考虑以下几种情景：

\begin{figure}
\centering
\resizebox{3in}{!}{\includegraphics{cartoons/r-2014-CPU-future-uniprocessor-uber-alles}}
\caption{单处理器至上}
\ContributedBy{fig:future:Uniprocessor \"Uber Alles}{Melissa Broussard}
\end{figure}

\begin{figure}
\centering
\resizebox{2.6in}{!}{\includegraphics{cartoons/r-2014-CPU-Future-Multithreaded-Mania}}
\caption{多线程狂热}
\ContributedBy{fig:future:Multithreaded Mania}{Melissa Broussard}
\end{figure}

\begin{figure}
\centering
\resizebox{2.5in}{!}{\includegraphics{cartoons/r-2014-CPU-Future-More-of-the-Same}}
\caption{一切照旧}
\ContributedBy{fig:future:More of the Same}{Melissa Broussard}
\end{figure}

\begin{figure}
\centering
\resizebox{3in}{!}{\includegraphics{cartoons/r-2014-CPU-Future-Crash-dummies}}
\caption{撞向内存墙的碰撞试验假人}
\ContributedBy{fig:future:Crash Dummies Slamming into the Memory Wall}{Melissa Broussard}
\end{figure}

\begin{figure}
\centering
\resizebox{3in}{!}{\includegraphics{cartoons/r-2021-CPU-future-astounding-accelerator}}
\caption{令人惊叹的加速器}
\ContributedBy{fig:future:Astounding Accelerators}{Melissa Broussard, remixed}
\end{figure}

\begin{enumerate}
\item	单处理器至上
	(\cref{fig:future:Uniprocessor \"Uber Alles})，
\item	多线程狂热
	(\cref{fig:future:Multithreaded Mania})，
\item	一切照旧
	(\cref{fig:future:More of the Same})，以及
\item	撞向内存墙的碰撞试验假人
	(\cref{fig:future:Crash Dummies Slamming into the Memory Wall})。
\item	令人惊叹的加速器
	(\cref{fig:future:Astounding Accelerators})。
\end{enumerate}

以下各节分别介绍每种情景。

\subsection{单处理器至上}
\label{sec:future:Uniprocessor \"Uber Alles}

正如2004年所说~\cite{PaulEdwardMcKenneyPhD}：

\begin{quote}
	在这一情景中，\IXaltr{Moore's-Law}{Moore's Law}带来的CPU时钟频率提升
	与水平扩展计算的持续进步相结合，使得SMP系统变得无关紧要。
	因此，这一情景被称为``单处理器至上''，
	字面意思是单处理器高于一切。

	这些单处理器系统将只受到指令开销的影响，因为\IXpl{memory barrier}、
	缓存抖动和竞争不会影响单CPU系统。
	在这一情景中，RCU仅对小众应用有用，例如
	与\IXacrpl{nmi}交互。
	目前还不清楚一个缺少RCU的操作系统是否会觉得有必要采用它，
	尽管已经实现了RCU的操作系统可能会继续使用。

	然而，多线程CPU的最新进展似乎表明
	这一情景是极不可能的。
\end{quote}

确实极不可能！
但更广泛的软件社区不愿接受他们需要拥抱并行性这一事实，
还需要一段时日，这个社区才会得出结论：
\IXaltr{Moore's-Law}{Moore's Law}带来的CPU时钟频率提升这顿``免费午餐''
确实已经彻底结束了。
永远不要忘记：
信念是一种情感，而不是理性技术思考的结果！

\subsection{多线程狂热}
\label{sec:future:Multithreaded Mania}

同样来自2004年~\cite{PaulEdwardMcKenneyPhD}：

\begin{quote}
	单处理器至上有一种不那么极端的变体，即具有硬件多线程的单处理器，
	事实上，多线程CPU现在已经是许多台式和笔记本电脑系统的标配。
	最激进的多线程CPU共享所有级别的缓存层次结构，
	从而消除了CPU间的\IXh{memory}{latency}，
	相应地也大大降低了传统同步机制的性能损失。
	然而，多线程CPU仍然会因竞争和内存屏障（memory barrier）导致的
	流水线停顿而产生开销。
	此外，由于所有硬件线程共享所有级别的缓存，
	单个硬件线程可用的缓存只是等效单线程CPU上缓存的一小部分，
	这可能会降低大量使用缓存的应用程序的性能。
	还存在一种可能性，即有限的可用缓存将导致基于RCU的算法
	因其 grace period 引起的额外内存消耗而遭受性能损失。
	调查这种可能性是未来的工作。

	然而，为了避免这种性能下降，许多多线程CPU和多CPU芯片
	至少在某些缓存级别上按每个硬件线程进行分区。
	这增加了每个硬件线程可用的缓存量，
	但重新引入了硬件线程间传递缓存行的内存延迟。
\end{quote}

最终，我们都知道这个故事是如何发展的：
单个芯片上有多个多线程核心，插在单个插槽中，
对较少的活跃线程数有不同程度的优化。
问题是，未来的共享内存系统是否总是能装入单个插槽中。

\subsection{一切照旧}
\label{sec:meas:More of the Same}

再次来自2004年~\cite{PaulEdwardMcKenneyPhD}：

\begin{quote}
	一切照旧情景假设，内存延迟比率将大致保持在今天的水平。

	这一情景实际上代表了一种变化，因为要保持不变，
	互连性能必须跟上\IXaltr{Moore's-Law}{Moore's Law}
	带来的核心CPU性能提升。
	在这一情景中，由于流水线停顿、内存延迟和竞争引起的开销
	仍然十分显著，RCU仍然保持着它今天所享有的高适用性。
\end{quote}

而变化正是\IXr{Moore's Law}仍在提供的不断提升的集成度。
但从长远来看，到底是什么？
在每个芯片上放置更多的CPU？
还是更多的I/O、缓存和内存？

服务器似乎在选择前者，而片上系统（SoC）的嵌入式系统则继续选择后者。

\subsection{撞向内存墙的碰撞试验假人}
\label{sec:future:Crash Dummies Slamming into the Memory Wall}

\begin{figure}
\centering
\epsfxsize=3in
\epsfbox{future/latencytrend}
% from Ph.D. thesis: related/latencytrend.eps
\caption{Sequent计算机每次本地内存引用的指令数}
\label{fig:future:Instructions per Local Memory Reference for Sequent Computers}
\end{figure}

\begin{figure}
\centering
\epsfxsize=3in
\epsfbox{future/be-lb-n4-rf-all}
% from Ph.D. thesis: an/plots/be-lb-n4-rf-all.eps
\caption{盈亏平衡点 vs.\@ $r$，$\lambda$ 较大，四个CPU}
\label{fig:future:Breakevens vs. r; lambda Large; Four CPUs}
\end{figure}

\begin{figure}
\centering
\epsfxsize=3in
\epsfbox{future/be-lw-n4-rf-all}
% from Ph.D. thesis: an/plots/be-lw-n4-rf-all.eps
\caption{盈亏平衡点 vs.\@ $r$，$\lambda$ 较小，四个CPU}
\label{fig:future:Breakevens vs. r; Worst-Case lambda; Four CPUs}
\end{figure}

2004年的又一段引言~\cite{PaulEdwardMcKenneyPhD}：

\begin{quote}
	如果
	\cref{fig:future:Instructions per Local Memory Reference for Sequent Computers}
	中所示的内存延迟趋势继续下去，那么内存延迟将继续相对于
	指令执行开销增长。
	像Linux这样大量使用RCU的系统将发现进一步使用RCU是有利可图的，
	如\cref{fig:future:Breakevens vs. r; lambda Large; Four CPUs}所示。
	从该图可以看出，如果RCU被大量使用，增加的内存延迟比率
	使RCU相对于其他同步机制获得越来越大的优势。
	相比之下，很少使用RCU的系统将需要越来越高的读密集度
	才能使RCU的使用有所回报，
	如\cref{fig:future:Breakevens vs. r; Worst-Case lambda; Four CPUs}所示。
	从该图可以看出，如果RCU使用较少，增加的内存延迟比率
	使RCU与其他同步机制相比处于越来越不利的地位。
	由于在高负载下Linux已被观察到每个\IX{grace period}
	超过1,600个回调~\cite{Sarma04c}，
	可以安全地说Linux属于前一类。
\end{quote}

一方面，这段话未能预见到缓存预热的问题，即RCU
在承受大量高强度更新写负载时可能遭受的性能损失，
大概是因为当时看来RCU不太可能用于这种工作负载。
实际上，\co{SLAB_TYPESAFE_BY_RCU}已在许多实例中投入使用，
在这些情况下否则会出现缓存预热问题，sequence locking也是如此。
另一方面，这段话也未能预见到RCU会被用于减少调度延迟或用于安全目的。

本书生成的大部分数据是在一个八插槽系统上收集的，
每个插槽有28个核心，每个核心有两个硬件线程，
总共448个硬件线程。
空闲系统的内存延迟小于一微秒，
与2004年同等规模系统的延迟不相上下。
有些人声称这些延迟接近一微秒仅仅是因为x86 CPU系列
相对较强的内存排序，但可能还需要一段时间
才能解决这一特定争论。

\subsection{令人惊叹的加速器}
\label{sec:future:Astounding Accelerators}

硬件加速器的潜力在2004年不像在2021年那样清晰，
因此本节没有引言。
然而，2020年11月的Top 500列表~\cite{Top500}包含大量加速器，
因此可以说本节是对现在而非未来的展望。
前面几节中的大多数也可以这样说。

硬件加速器正被用于许多其他用途，包括
加密、压缩和机器学习。

简而言之，还是要小心进行预测，包括本章其余部分的预测。
